% Part: first-order-logic
% Chapter: introduction
% Section: semantic-notions

\documentclass[../../../include/open-logic-section]{subfiles}

\begin{document}

\olfileid[hi]{fol}{int}{sem}

\olsection{अर्थवैज्ञानिक धारणाएँ}

ऊपर हमने कहा कि $\Sat{M}{!A}[s]$ सत्य है या नहीं, इस पर विचार करते समय
सुविधा के लिए $s$ को सभी !!{variable}s के मान देने देते हैं, किंतु उपयोग
केवल उन मानों का होता है जो वह~$!A$ में आने वाले !!{variable}s को देता है।
वास्तव में केवल~$!A$ के \emph{मुक्त} चरों के मान महत्त्व रखते हैं। सावधानी
के कारण हम इस तथ्य को सिद्ध भी करेंगे। चूँकि !!{sentence}s में कोई मुक्त चर
नहीं होता, इसलिए !!a{structure} में वे संतुष्ट हैं या नहीं, इसमें~$s$ का कोई
महत्त्व नहीं। अतः जब $!A$ कोई !!a{sentence} हो, तो $\Sat{M}{!A}$ का अर्थ
``$\Sat{M}{!A}[s]$, सभी~$s$ के लिए'' परिभाषित किया जा सकता है। यह ठीक तभी
सत्य होता है जब $\Sat{M}{!A}[s]$ कम-से-कम एक~$s$ के लिए सत्य हो।
!!{formula}s के लिए संतुष्टि की उपयोगी परिभाषा पाने हेतु !!{variable}
निर्धारण आवश्यक हैं, किंतु !!{sentence}s के लिए संतुष्टि !!{variable}
निर्धारणों से स्वतंत्र है।

``$\Sat{M}{!A}$'' की परिभाषा मिल जाने पर प्रथम-क्रम तर्क के !!{sentence}s
के संदर्भ में हम जानते हैं कि ``स्थिति'' और ``में सत्य'' का क्या अर्थ है।
!!{sentence}s के लिए $\Sat{M}{!A}$ की परिभाषा के आधार पर हम वैधता, निष्पत्ति
और संतुष्टनीयता की मूल अर्थवैज्ञानिक धारणाएँ परिभाषित कर सकते हैं। कोई वाक्य
तब वैध, $\Entails !A$, है जब प्रत्येक !!{structure} उसे संतुष्ट करती है। वह
!!{sentence}s के किसी समुच्चय से निष्पन्न, $\Gamma \Entails !A$, तब होता है
जब~$\Gamma$ के सभी !!{sentence}s को संतुष्ट करने वाली प्रत्येक
!!{structure},~$!A$ को भी संतुष्ट करती है। और !!{sentence}s का कोई समुच्चय तब
संतुष्टनीय है जब कोई !!{structure} उसके सभी !!{sentence}s को एक साथ संतुष्ट
करती हो।

चूँकि !!{formula}s का वर्ग आगमनात्मक रूप से परिभाषित है और संतुष्टि भी !!{formula}s
की संरचना पर आगमन द्वारा परिभाषित है, इसलिए अर्थविज्ञान के गुण सिद्ध करने और
परिभाषित अर्थवैज्ञानिक धारणाओं को परस्पर जोड़ने के लिए हम आगमन का उपयोग कर
सकते हैं। इनमें से कुछ गुणों को हम संकलित और सिद्ध करेंगे: आंशिक रूप से इसलिए
कि वे अपने-आप में रोचक हैं, किंतु मुख्यतः इसलिए कि अर्थविज्ञान और
!!{derivation} प्रणालियों के संबंध की आगे जाँच करते समय उनमें से अनेक उपयोगी
होंगे। इसके लिए कुछ ऐसी वाक्यविन्यासी धारणाएँ और संक्रियाएँ भी सटीक रूप से,
अर्थात् आगमन द्वारा, परिभाषित करनी होंगी जिनका अभी उल्लेख नहीं हुआ।

\end{document}
